\section{Introduction}
\label{sec:intro}

Traffic congestion at signalized intersections costs travel time, fuel and safety in cities worldwide. Most intersections still run fixed-time plans, which repeat the same cycle regardless of traffic, or actuated control, which stretches or cuts a green phase when loop detectors report arriving vehicles~\cite{michailidis2025review}. Deep reinforcement learning (DRL) offers a data-driven alternative: an agent at each intersection observes the queues on its approaches and chooses the next green phase, learning a policy that reduces delay~\cite{michailidis2025review,mei2024libsignal}.

Two gaps make it hard to tell how such controllers would perform in practice. First, evaluation rarely separates training and test demand. A scenario typically ships one hour of calibrated traffic, and many studies train and report on that same hour, so a controller can perform well by memorizing one demand pattern~\cite{ault2021resco,kwesiga2026robustness}. Second, algorithm comparisons are rarely controlled: learners are reported with different observations, rewards, yellow times and numbers of seeds, so differences can come from the setup rather than the algorithm~\cite{mei2024libsignal,agarwal2021precipice}.

This work is useful to traffic agencies and researchers who must choose a DRL controller before deploying it: it shows which widely used algorithm keeps its benefit when traffic differs from the training data, how often each one fails, and what the failures look like, and it provides a reusable protocol for such tests. It is also the first stage of a larger study on resource-aware edge--cloud intelligent transportation systems. We compare DQN, Double DQN, PPO and A2C because they represent the two main DRL families used for signal control: value-based off-policy learning (DQN, and Double DQN, which corrects DQN's over-estimation of action values) and on-policy actor-critic learning (A2C, and PPO, which limits the size of each policy update). Independent DQN and PPO learners were also the strongest methods on the RESCO benchmark~\cite{ault2021resco}, and all four have mature, directly comparable implementations~\cite{raffin2021sb3}. As real-city scenarios we use cologne1, ingolstadt1 and cologne3 from RESCO~\cite{ault2021resco}, which provide real road geometry and calibrated peak-hour demand from the TAPASCologne and InTAS traffic models.

The contributions of this study are: (i) a generalization-focused protocol that trains on $0.8$--$1.2\times$ variants of the real peak hour and tests only on unseen $0.7\times$ and $1.3\times$ demand and on the original hour; (ii) a controlled comparison of DQN, Double DQN, PPO and A2C against fixed-time, tuned actuated and random control with identical observations, rewards, yellow times and interaction budgets, and five seeds each; (iii) a failure analysis showing that rare gridlock episodes, rather than average performance, separate the learners; and (iv) an ablation of the two training components, demand randomization and policy sharing, and of valid-action masking, with per-intersection results that locate the failures. The generalization-focused evaluation of DRL signal controllers on demand levels not encountered during training is considered the main novelty.
